\section{Homework Week 7}

\begin{exercise}
    Let \(V\) be an \(A\)-module.
    \begin{enumerate}
        \renewcommand{\labelenumi}{(\alph{enumi})}
        \item Show that an endomorphism
              \[
                  e \colon V \to V
              \]
              is a projection onto an \(A\)-submodule \(W\) if and only if \(e\) is idempotent as an element of \(\operatorname{End}_A(V)\).

        \item Show that direct sum decompositions
              \[
                  V = W_1 \oplus W_2
              \]
              as \(A\)-modules are in bijection with idempotent decompositions
              \[
                  1 = e + f
              \]
              in \(\operatorname{End}_A(V)\).

        \item Show that \(e \in \operatorname{End}_A(V)\) is primitive if and only if \(e(V)\) is indecomposable.

        \item Suppose that \(V\) is semisimple with finitely many simple summands and let
              \[
                  e_1, e_2 \in \operatorname{End}_A(V)
              \]
              be idempotents. Show that
              \[
                  e_1(V) \cong e_2(V)
              \]
              as \(A\)-modules if and only if \(e_1\) and \(e_2\) are conjugate by an invertible element of \(\operatorname{End}_A(V)\).

        \item Let \(k\) be a field. Show that all primitive idempotent elements in \(M_n(k)\) are conjugate under the action of the unit group \(GL_n(k)\), which consists of all the invertible matrices in \(M_n(k)\).
    \end{enumerate}
\end{exercise}

\begin{proof}
    \begin{enumerate}
        \renewcommand{\labelenumi}{(\alph{enumi})}

        \item Suppose first that \(e\colon V\to V\) is a projection onto an \(A\)-submodule \(W\).
              Then \(e(V)=W\) and \(e(w)=w\) for every \(w\in W\). Hence for any \(v\in V\),
              \[
                  e^2(v)=e(e(v))=e(v),
              \]
              because \(e(v)\in W\). Thus \(e^2=e\), so \(e\) is idempotent.

              Conversely, suppose that \(e^2=e\). Set
              \[
                  W:=e(V).
              \]
              Since \(e\) is \(A\)-linear, \(W\) is an \(A\)-submodule of \(V\). Moreover, if \(w=e(v)\in W\), then
              \[
                  e(w)=e^2(v)=e(v)=w.
              \]
              Thus \(e\) fixes \(W\) pointwise and maps \(V\) onto \(W\), so \(e\) is a projection onto \(W\).

              In fact, one also has
              \[
                  V=e(V)\oplus \ker(e),
              \]
              since for any \(v\in V\),
              \[
                  v=e(v)+(v-e(v)),
                  \qquad e(v-e(v))=e(v)-e^2(v)=0,
              \]
              and if \(x\in e(V)\cap \ker(e)\), then \(x=e(v)\) for some \(v\), hence
              \[
                  x=e(x)=0.
              \]

        \item Suppose first that
              \[
                  V=W_1\oplus W_2.
              \]
              Define
              \[
                  e(w_1+w_2)=w_1,
                  \qquad
                  f(w_1+w_2)=w_2
                  \qquad (w_i\in W_i).
              \]
              Then \(e,f\in \operatorname{End}_A(V)\), and \(e\) and \(f\) are the projections onto
              \(W_1\) and \(W_2\), respectively. By part (a), both are idempotent, and clearly
              \[
                  e+f=1.
              \]
              Also,
              \[
                  ef=fe=0.
              \]

              Conversely, suppose
              \[
                  1=e+f
              \]
              with \(e^2=e\) and \(f^2=f\). Then
              \[
                  e=e(e+f)=e+ef,
              \]
              so \(ef=0\). Similarly,
              \[
                  f=(e+f)f=ef+f,
              \]
              so again \(ef=0\), and likewise
              \[
                  fe=0.
              \]
              By part (a), \(e\) and \(f\) are projections onto \(e(V)\) and \(f(V)\), respectively.
              For every \(v\in V\),
              \[
                  v=(e+f)(v)=e(v)+f(v),
              \]
              so
              \[
                  V=e(V)+f(V).
              \]
              If \(x\in e(V)\cap f(V)\), say \(x=e(v)=f(w)\), then
              \[
                  x=e(x)=ef(w)=0.
              \]
              Hence
              \[
                  V=e(V)\oplus f(V).
              \]

              These two constructions are inverse to each other, so direct sum decompositions
              \[
                  V=W_1\oplus W_2
              \]
              are in bijection with idempotent decompositions
              \[
                  1=e+f
              \]
              in \(\operatorname{End}_A(V)\).

        \item We show that \(e\) is primitive if and only if \(e(V)\) is indecomposable.

              First suppose that \(e\) is not primitive. Then there exist nonzero idempotents
              \(e_1,e_2\in \operatorname{End}_A(V)\) such that
              \[
                  e=e_1+e_2,
                  \qquad
                  e_1e_2=e_2e_1=0.
              \]
              Since
              \[
                  ee_i=(e_1+e_2)e_i=e_i,
              \]
              we have \(e_i(V)\subseteq e(V)\) for \(i=1,2\). Moreover, for any \(v\in V\),
              \[
                  e(v)=e_1(v)+e_2(v),
              \]
              so
              \[
                  e(V)=e_1(V)+e_2(V).
              \]
              If \(x\in e_1(V)\cap e_2(V)\), say \(x=e_1(v)=e_2(w)\), then
              \[
                  x=e_1(x)=e_1e_2(w)=0.
              \]
              Hence
              \[
                  e(V)=e_1(V)\oplus e_2(V).
              \]
              Since \(e_1\) and \(e_2\) are nonzero, both \(e_1(V)\) and \(e_2(V)\) are nonzero.
              Thus \(e(V)\) is decomposable.

              Conversely, suppose that \(e(V)\) is decomposable:
              \[
                  e(V)=U_1\oplus U_2
              \]
              with \(U_1,U_2\neq 0\). Let
              \[
                  p_i\colon e(V)\to e(V)
              \]
              be the projection onto \(U_i\) (\(i=1,2\)). Define
              \[
                  \widetilde e_i:=p_i\circ e \in \operatorname{End}_A(V).
              \]
              Since \(e\) acts as the identity on \(e(V)\), we have \(e\circ p_i=p_i\), and therefore
              \[
                  \widetilde e_i^2=p_i e p_i e=p_i p_i e=\widetilde e_i.
              \]
              Also,
              \[
                  \widetilde e_1\widetilde e_2=p_1 e p_2 e=p_1p_2 e=0,
                  \qquad
                  \widetilde e_2\widetilde e_1=0,
              \]
              and
              \[
                  \widetilde e_1+\widetilde e_2=(p_1+p_2)e=e.
              \]
              Since \(U_i\neq 0\), each \(\widetilde e_i\neq 0\). Thus \(e\) is not primitive.

              Hence \(e\) is primitive if and only if \(e(V)\) is indecomposable.

        \item Assume first that \(e_2=u e_1 u^{-1}\) for some invertible
              \(u\in \operatorname{End}_A(V)\). Then for any \(x=e_1(v)\in e_1(V)\),
              \[
                  u(x)=u e_1(v)=e_2u(v)\in e_2(V),
              \]
              so \(u(e_1(V))\subseteq e_2(V)\). Conversely, if \(y=e_2(w)\in e_2(V)\), then
              \[
                  y=u e_1 u^{-1}(w)\in u(e_1(V)).
              \]
              Hence
              \[
                  u(e_1(V))=e_2(V),
              \]
              and therefore the restriction
              \[
                  u|_{e_1(V)}\colon e_1(V)\to e_2(V)
              \]
              is an isomorphism of \(A\)-modules. So
              \[
                  e_1(V)\cong e_2(V).
              \]

              Conversely, suppose
              \[
                  e_1(V)\cong e_2(V).
              \]
              Since \(V\) is semisimple with finitely many simple summands, we may write
              \[
                  V\cong \bigoplus_{j=1}^r S_j^{\,n_j},
              \]
              where \(S_1,\dots,S_r\) are pairwise nonisomorphic simple \(A\)-modules.
              Because \(e_i(V)\) is a submodule of a semisimple module, it is semisimple, so
              \[
                  e_i(V)\cong \bigoplus_{j=1}^r S_j^{\,m_j}
              \]
              for some integers \(m_j\), and the \(m_j\) are the same for \(i=1,2\) since
              \(e_1(V)\cong e_2(V)\).

              By part (a),
              \[
                  V=e_i(V)\oplus \ker(e_i).
              \]
              Therefore
              \[
                  \ker(e_i)\cong \bigoplus_{j=1}^r S_j^{\,n_j-m_j},
              \]
              and hence
              \[
                  \ker(e_1)\cong \ker(e_2).
              \]

              Choose isomorphisms
              \[
                  \alpha\colon e_1(V)\xrightarrow{\sim} e_2(V),
                  \qquad
                  \beta\colon \ker(e_1)\xrightarrow{\sim} \ker(e_2).
              \]
              Using the decompositions
              \[
                  V=e_1(V)\oplus \ker(e_1)
                  \qquad\text{and}\qquad
                  V=e_2(V)\oplus \ker(e_2),
              \]
              define
              \[
                  u\colon V\to V,
                  \qquad
                  u(x+y)=\alpha(x)+\beta(y)
                  \quad
                  \bigl(x\in e_1(V),\ y\in \ker(e_1)\bigr).
              \]
              Then \(u\) is an invertible \(A\)-endomorphism of \(V\). For \(x\in e_1(V)\) and
              \(y\in \ker(e_1)\),
              \[
                  ue_1(x+y)=u(x)=\alpha(x),
              \]
              while
              \[
                  e_2u(x+y)=e_2\bigl(\alpha(x)+\beta(y)\bigr)=\alpha(x),
              \]
              because \(e_2\) is the identity on \(e_2(V)\) and zero on \(\ker(e_2)\).
              Hence
              \[
                  ue_1=e_2u,
              \]
              so
              \[
                  e_2=ue_1u^{-1}.
              \]
              Thus \(e_1\) and \(e_2\) are conjugate by an invertible element of
              \(\operatorname{End}_A(V)\).

        \item Let \(V=k^n\). Then
              \[
                  M_n(k)\cong \operatorname{End}_k(V).
              \]
              Let \(e\in M_n(k)\) be a primitive idempotent. By part (c), \(e(V)\) is indecomposable
              as a \(k\)-vector space. But a nonzero \(k\)-vector space is indecomposable if and only if
              it has dimension \(1\): indeed, if \(\dim_k W\ge 2\), then any nonzero proper subspace has
              a complement, so \(W\) is decomposable.

              Therefore
              \[
                  \dim_k e(V)=1.
              \]
              Since \(e\) is idempotent, part (a) gives
              \[
                  V=e(V)\oplus \ker(e).
              \]
              Choose a basis \(v_1,\dots,v_n\) of \(V\) such that \(v_1\) spans \(e(V)\) and
              \(v_2,\dots,v_n\) is a basis of \(\ker(e)\). Then
              \[
                  e(v_1)=v_1,
                  \qquad
                  e(v_i)=0 \quad (i\ge 2).
              \]
              Hence the matrix of \(e\) with respect to this basis is
              \[
                  E_{11}=
                  \begin{pmatrix}
                      1      & 0      & \cdots & 0      \\
                      0      & 0      & \cdots & 0      \\
                      \vdots & \vdots & \ddots & \vdots \\
                      0      & 0      & \cdots & 0
                  \end{pmatrix}.
              \]
              Therefore \(e\) is conjugate to \(E_{11}\) by some element of \(GL_n(k)\). It follows that
              any two primitive idempotents in \(M_n(k)\) are conjugate to each other.
    \end{enumerate}
\end{proof}

\begin{exercise}
    Prove: No idempotent \(e\) of \(A\) is contained in \(J(A)\). [Hint: \(e(1-e)=0\).]
\end{exercise}

\begin{proof}
    Assume that \(e\in J(A)\) and that \(e\) is idempotent, so
    \[
        e^2=e.
    \]
    Since \(e\in J(A)\), a standard characterization of the Jacobson radical implies that
    \[
        1-e
    \]
    is a unit in \(A\).

    On the other hand,
    \[
        e(1-e)=e-e^2=0.
    \]
    Multiplying this equality on the right by \((1-e)^{-1}\), we obtain
    \[
        e=0.
    \]

    Therefore any idempotent contained in \(J(A)\) must be zero. In particular, no nonzero idempotent of \(A\) lies in \(J(A)\).
\end{proof}

\begin{exercise}
    Prove: Let \(A\) be an Artinian ring with Artinian center \(Z(A)\). Then
    \[
        J(Z(A)) = J(A) \cap Z(A).
    \]
\end{exercise}

\begin{proof}
    Let
    \[
        B:=Z(A).
    \]
    We prove the two inclusions separately.

    First, we show that
    \[
        J(B)\subseteq J(A)\cap B.
    \]
    Take \(z\in J(B)\). Since \(B\) is Artinian, its Jacobson radical is nilpotent, so there exists
    \(m\ge 1\) such that
    \[
        J(B)^m=0.
    \]
    In particular,
    \[
        z^m=0.
    \]
    Now let \(x\in A\). Because \(z\in Z(A)\), we have \(xz=zx\), hence
    \[
        (xz)^m=x^m z^m=0.
    \]
    So \(xz\) is nilpotent, and therefore
    \[
        1-xz
    \]
    is invertible in \(A\), with inverse
    \[
        1+xz+\cdots +(xz)^{m-1}.
    \]
    By the standard characterization of the Jacobson radical, this implies
    \[
        z\in J(A).
    \]
    Since also \(z\in B\), we get
    \[
        z\in J(A)\cap B.
    \]
    Thus
    \[
        J(B)\subseteq J(A)\cap B.
    \]

    Next, we show that
    \[
        J(A)\cap B\subseteq J(B).
    \]
    Take \(z\in J(A)\cap B\). Since \(A\) is Artinian, \(J(A)\) is nilpotent, so there exists
    \(n\ge 1\) such that
    \[
        J(A)^n=0.
    \]
    Hence
    \[
        z^n=0.
    \]
    Now let \(c\in B\). Since \(B\) is commutative, we have
    \[
        (cz)^n=c^n z^n=0.
    \]
    Therefore \(cz\) is nilpotent in \(B\), and so
    \[
        1-cz
    \]
    is invertible in \(B\), with inverse
    \[
        1+cz+\cdots +(cz)^{n-1}.
    \]
    Again by the standard characterization of the Jacobson radical, this shows that
    \[
        z\in J(B).
    \]
    Hence
    \[
        J(A)\cap B\subseteq J(B).
    \]

    Combining the two inclusions, we conclude that
    \[
        J(Z(A))=J(B)=J(A)\cap B.
    \]
\end{proof}

\begin{exercise}
    Prove: Let
    \[
        A = A_1 \oplus \cdots \oplus A_n
    \]
    be a direct sum of \(2\)-sided ideals. Then
    \begin{enumerate}
        \renewcommand{\labelenumi}{(\roman{enumi})}
        \item
              \[
                  Z(A) = Z(A_1) \oplus Z(A_2) \oplus \cdots \oplus Z(A_n).
              \]
        \item
              \[
                  J(A) = J(A_1) \oplus J(A_2) \oplus \cdots \oplus J(A_n).
              \]
    \end{enumerate}
\end{exercise}

\begin{proof}
    For \(i\neq j\), since \(A_i\) and \(A_j\) are two-sided ideals and the sum is direct,
    \[
        A_iA_j\subseteq A_i\cap A_j=0.
    \]
    Write
    \[
        1_A=e_1+\cdots+e_n,\qquad e_i\in A_i.
    \]
    Then \(e_i\) is the identity element of \(A_i\). Indeed, if \(a_i\in A_i\), then
    all mixed products vanish, so
    \[
        a_i=1_Aa_i=e_i a_i,\qquad
        a_i=a_i1_A=a_i e_i.
    \]
    Thus each \(A_i\) is a unital ring, with identity \(e_i\). Moreover multiplication in
    \(A\) is coordinatewise with respect to the decomposition
    \[
        A=A_1\oplus\cdots\oplus A_n.
    \]
    Hence
    \[
        \Phi\colon A\longrightarrow A_1\times\cdots\times A_n,
        \qquad
        a_1+\cdots+a_n\longmapsto (a_1,\dots,a_n)
    \]
    is an isomorphism of rings. Since \(n\) is finite, we identify
    \[
        A \cong A_1\times\cdots\times A_n.
    \]

    \begin{enumerate}
        \renewcommand{\labelenumi}{(\roman{enumi})}

        \item First we prove
              \[
                  Z(A)\subseteq Z(A_1)\oplus\cdots\oplus Z(A_n).
              \]
              Let \(z\in Z(A)\), and write
              \[
                  z=z_1+\cdots+z_n,\qquad z_i\in A_i.
              \]
              Fix \(i\), and let \(x_i\in A_i\). Since \(z\) is central in \(A\),
              \[
                  zx_i=x_i z.
              \]
              But all mixed products vanish, so
              \[
                  zx_i=z_i x_i,\qquad x_i z=x_i z_i.
              \]
              Hence \(z_i x_i=x_i z_i\) for every \(x_i\in A_i\), and therefore
              \[
                  z_i\in Z(A_i).
              \]
              Since this holds for every \(i\), we get
              \[
                  z\in Z(A_1)\oplus\cdots\oplus Z(A_n).
              \]

              Conversely, suppose
              \[
                  z=z_1+\cdots+z_n,\qquad z_i\in Z(A_i).
              \]
              Let
              \[
                  a=a_1+\cdots+a_n,\qquad a_i\in A_i.
              \]
              Again using \(A_iA_j=0\) for \(i\neq j\), we have
              \[
                  za=\sum_{i=1}^n z_i a_i,
                  \qquad
                  az=\sum_{i=1}^n a_i z_i.
              \]
              Since \(z_i\in Z(A_i)\), each \(z_i a_i=a_i z_i\). Hence \(za=az\), so
              \(z\in Z(A)\). Therefore
              \[
                  Z(A)=Z(A_1)\oplus\cdots\oplus Z(A_n).
              \]

        \item First we prove
              \[
                  J(A_1)\oplus\cdots\oplus J(A_n)\subseteq J(A).
              \]
              Let
              \[
                  x=x_1+\cdots+x_n,\qquad x_i\in J(A_i).
              \]
              Take any
              \[
                  a=a_1+\cdots+a_n\in A,\qquad a_i\in A_i.
              \]
              Since \(J(A_i)\) is an ideal of \(A_i\), we have \(a_i x_i\in J(A_i)\). Hence
              \(e_i-a_i x_i\) is invertible in \(A_i\). Choose \(u_i\in A_i\) such that
              \[
                  u_i(e_i-a_i x_i)=e_i.
              \]
              Put \(u=u_1+\cdots+u_n\). Since multiplication is coordinatewise,
              \[
                  u(1_A-ax)
                  =
                  \sum_{i=1}^n u_i(e_i-a_i x_i)
                  =
                  \sum_{i=1}^n e_i
                  =
                  1_A.
              \]
              Therefore \(1_A-ax\) has a left inverse for every \(a\in A\). By the Jacobson
              radical criterion,
              \[
                  x\in J(A).
              \]
              Thus
              \[
                  J(A_1)\oplus\cdots\oplus J(A_n)\subseteq J(A).
              \]

              Conversely, we prove
              \[
                  J(A)\subseteq J(A_1)\oplus\cdots\oplus J(A_n).
              \]
              Let
              \[
                  x=x_1+\cdots+x_n\in J(A),\qquad x_i\in A_i.
              \]
              Fix \(i\), and let \(a_i\in A_i\). Since \(x\in J(A)\), the element
              \[
                  1_A-a_i x=1_A-a_i x_i
              \]
              has a left inverse in \(A\). Thus there exists \(v=v_1+\cdots+v_n\in A\), with
              \(v_j\in A_j\), such that
              \[
                  v(1_A-a_i x_i)=1_A.
              \]
              Taking the \(A_i\)-component of this equality gives
              \[
                  v_i(e_i-a_i x_i)=e_i.
              \]
              Hence \(e_i-a_i x_i\) has a left inverse in \(A_i\) for every \(a_i\in A_i\). By the
              Jacobson radical criterion applied inside \(A_i\),
              \[
                  x_i\in J(A_i).
              \]
              Since \(i\) was arbitrary,
              \[
                  x\in J(A_1)\oplus\cdots\oplus J(A_n).
              \]
              Combining the two inclusions, we conclude that
              \[
                  J(A)=J(A_1)\oplus\cdots\oplus J(A_n).
              \]
    \end{enumerate}
\end{proof}

\begin{exercise}
    Let
    \[
        \phi \colon U \to V
    \]
    be a homomorphism of \(A\)-modules. Show that
    \[
        \phi(\operatorname{Soc} U) \subseteq \operatorname{Soc} V,
    \]
    and that if \(\phi\) is an isomorphism then \(\phi\) restricts to an isomorphism
    \[
        \operatorname{Soc} U \to \operatorname{Soc} V.
    \]
\end{exercise}

\begin{proof}
    Recall that
    \[
        \operatorname{Soc}(U)=\sum \{\,S\leq U \mid S \text{ is a simple } A\text{-submodule}\,\}.
    \]

    We first show that
    \[
        \phi(\operatorname{Soc} U)\subseteq \operatorname{Soc} V.
    \]
    Let \(S\leq U\) be a simple submodule. Then \(\phi(S)\) is a submodule of \(V\), and since
    \(S\) is simple, the restriction
    \[
        \phi|_S \colon S\to \phi(S)
    \]
    is either the zero map or an isomorphism. Hence \(\phi(S)\) is either \(0\) or a simple
    submodule of \(V\). In either case,
    \[
        \phi(S)\subseteq \operatorname{Soc} V.
    \]
    Since \(\operatorname{Soc}(U)\) is the sum of all simple submodules \(S\le U\), we get
    \[
        \phi(\operatorname{Soc} U)
        =
        \sum_{S\text{ simple }\le U}\phi(S)
        \subseteq \operatorname{Soc} V.
    \]

    Now assume that \(\phi\) is an isomorphism. Then \(\phi^{-1}\colon V\to U\) is also an
    \(A\)-module homomorphism. Applying the first part to \(\phi^{-1}\), we obtain
    \[
        \phi^{-1}(\operatorname{Soc} V)\subseteq \operatorname{Soc} U.
    \]
    Applying \(\phi\) to both sides gives
    \[
        \operatorname{Soc} V \subseteq \phi(\operatorname{Soc} U).
    \]
    Together with the already proved inclusion
    \[
        \phi(\operatorname{Soc} U)\subseteq \operatorname{Soc} V,
    \]
    we conclude that
    \[
        \phi(\operatorname{Soc} U)=\operatorname{Soc} V.
    \]

    Therefore the restriction
    \[
        \phi|_{\operatorname{Soc} U}\colon \operatorname{Soc} U\to \operatorname{Soc} V
    \]
    is surjective. It is also injective because \(\phi\) itself is injective. Hence
    \[
        \phi|_{\operatorname{Soc} U}\colon \operatorname{Soc} U \xrightarrow{\sim} \operatorname{Soc} V
    \]
    is an isomorphism of \(A\)-modules.
\end{proof}

\begin{exercise}
    If
    \[
        \varphi \colon A \to A'
    \]
    is an epimorphism of rings, then
    \[
        \varphi(J(A)) \subseteq J(A').
    \]
\end{exercise}

\begin{proof}
    Let \(a\in J(A)\). We show that \(\varphi(a)\in J(A')\).

    We use the standard characterization of the Jacobson radical:
    \[
        x\in J(R)
        \quad\Longleftrightarrow\quad
        \text{for every } r\in R,\ \ 1-rx \text{ has a left inverse in } R.
    \]

    So let \(y'\in A'\). Since \(\varphi\colon A\to A'\) is surjective, there exists \(y\in A\)
    such that
    \[
        \varphi(y)=y'.
    \]
    Because \(a\in J(A)\), the element
    \[
        1-ya
    \]
    has a left inverse in \(A\). Thus there exists \(b\in A\) such that
    \[
        b(1-ya)=1.
    \]
    Apply \(\varphi\) to this equality. Since \(\varphi\) is a ring homomorphism, we get
    \[
        \varphi(b)\,\varphi(1-ya)=\varphi(1)=1,
    \]
    that is,
    \[
        \varphi(b)\bigl(1-\varphi(y)\varphi(a)\bigr)=1.
    \]
    Using \(\varphi(y)=y'\), this becomes
    \[
        \varphi(b)(1-y'\varphi(a))=1.
    \]
    Hence \(1-y'\varphi(a)\) has a left inverse in \(A'\).

    Since \(y'\in A'\) was arbitrary, the above characterization implies that
    \[
        \varphi(a)\in J(A').
    \]
    Therefore
    \[
        \varphi(J(A))\subseteq J(A').
    \]
\end{proof}

